
\documentclass{article}

\usepackage[czech]{babel} 
\usepackage[utf8]{inputenc}
\usepackage[T1]{fontenc}
%\usepackage[IL2]{fontenc} % Different type of font for wierdos
\usepackage{graphicx} % Required for the inclusion of images
\usepackage{natbib} % Required to change bibliography style to APA
\usepackage{amsmath} % Required for some math elements 
\usepackage{amsfonts}
\usepackage{amssymb}
\usepackage{booktabs}
\usepackage{epstopdf}
\usepackage{multicol}
\usepackage{color}
\usepackage{float}
\usepackage[total={15.5cm,23.5cm}, top=2.5cm, left=3cm, includefoot]{geometry}
\usepackage[colorlinks=true,linkcolor = black, urlcolor = black, citecolor = black]{hyperref}
\usepackage{enumitem}
\usepackage{subfig}
\usepackage[toc,page]{appendix}
\usepackage{hyperref} 
\usepackage{matlab-prettifier}
\usepackage{datetime}
\usepackage{mathtools}
\usepackage[thinc]{esdiff}
\newcommand{\myparagraph}[1]{\paragraph{#1}\mbox{}\\}

\begin{document}


\begin{titlepage}

    \centering

    {\scshape\LARGE Západočeská univerzita\par}
    {\scshape\Large Fakulta aplikovaných věd \par}
    {\scshape\Large Katerdra Kybernetiky \par}
    {\begin{center}
            \includegraphics[width=0.7\textwidth]{../pic/fav.jpg}
        \end{center}}

    {\huge\bfseries Dokumentace distribuci projeku ITE-YELLOW\par} % TODO: lepsi nazev

    \vspace{2cm}

    {\Large\itshape Martin Hamar, Radek Kaupe, Samuel Kokoška\par}

    \vfill

    \vspace{1cm}

    \hskip -1 cm{KKY/BSOI} \hfill {Datum: \today }



\end{titlepage}

\section{Úvod}
V rámci předmětu KKY/BSOI jsme si vytvořili a nastavili virtuální stroj s operačním systémem Linux. 
Tento stroj je momentálně používán k provozu applikace vyvíjené v rámci předmětu KKY/ITE. \href{http://147.228.173.162:8881}{Zde} je na aplikaci a \href{https://github.com/RadekKaupe/ITE-Yellow-2024}{zde} je odkaz na Githubový repozitář, kde jsou k dipozici veškeré zdojové kódy aplikace jak části ITE, tak části BSOI. 
Tento referát by měl popsat virtuální stroj, jeho nastavení, činnosti a realizaci způsobů přihlášení do již zmíněné aplikace.   
\section{Řešení}

\subsection{Vlastnosti virtuálního stroje}
Při zakládání stroje (sulis162, 147.228.173.162) byly zvoleny tyto parametry; 1 procesor se 2 vlákny, 2 GB operační paměťi a 20 GB paměti pevného disku. Na stroji běží operační systém Debian 11.6.
\section{Uživatelské účty}
Na virtuálním stroji jsou tři uživatelské účty - každý pro jednoho člena týmu. Tyto účty mohou používat sudo a mohou se na stroj přihlásit přes ssh bez hesla díky jejich veřejným ssh klíčům.
Přímé přihlášení na root pomocí ssh je vypnuté.
\subsection{Monitoring}
Dva členové týmu využil pro monitoring služeb UptimeRobot (\href{https://uptimerobot.com/}{https://uptimerobot.com/}), kde byl nastaven ping na IP adresu virtuálního stroje každých 5 minut. V výpadku nás služba upozorní emailem.

\subsection{Zálohování}
Zálohování databáze bylo implementováno pomocí nástroje Borgbackup, specificky funkcí prune pro automatizaci celého procesu k ukládání každodenní zálohy na účty skokoska a radek, v případě poruchy či potřeby vrátit databázi do předchozího stavu tedy existují dvě zálohy.

Na každém účtě je složka `backup' obsahující zašifrovanou zálohu pod defaultním heslem `gobackImadeamistak'.
\subsection{Šifrování}
Certifikáty  pro šifrování pomocí TLS (Transport Layer Security) byly vystaveny od autority `Let's encrypt' a byl použit šifrovací klíč ECC. Certifikát je platný 90 dní a má celkem tři části:
\begin{enumerate}
    \item key.pem - obsahuje soukromý klíč, který je pužit k dešifrování dat
    \item cert.pem - obsahuje certifikát serve  ru
    \item fullchain.pem - řetězec důvěry
\end{enumerate}.
\subsection{Zabezpečení virtuálního stroje}
Samotnáwebová aplikace je přístupná pouze ze sítě ZČU a to na portu 442, zde také v našem případě probíhá HTTPS komunikace.
Poté bylo třeba otevřít port 80 (HTTP) kvůli vystavení TLS certifikátu.
Žádné jiné porty otevřené nejsou.
\subsection{Provoz distribuované aplikace}
Součásí aplikace je PostgreSQL databáze a dva python skripty, jeden z nich řeší MQTT subscribera, abychom mohli data příjímat, validovat a ukládat a druhý tvoří backend naší webové aplikace.
K založení PostgreSQL datbáze, byl také vytovřen python skript, který stačí spustit pouze jedno, aby vytvořil tabulky. Skripty na subscribera a backend musí běžet neustále, aby aplikace mohla být v provozu.
Tyto skript běží současně pomocí tmuxu. Tmux je `Terminal Multiplexer', který umožňuje rozdělit obrazovku terminálu na více panelů a vytváření `sezení' (sessions), které mohou běžet i na pozadí, poté co se člověk odpojí.
Subscriber a backend spolu `komunikují' pouze přes databázi, ke které mají oba skripty přístup.

ZAJISTI AUTOMATICKY SPUSTENI PO RESTARTu

Momentálně logování výstupu je k vidění pouze v rámci tmux session, kde jsou nastavené výpisy za běhu obou skriptů, aby mohlo dojít ke kontrole kdykoliv některým členů z týmu. 
Nasazení verzí se může provést pomocí programu WinSCP, kterým se nové soubory nahrají na virtuální stroj.
\subsection{Přihlášení a biometrická autentizace}
Při implementaci jsem využil učiteli dodaných zdrojů, ve formě skriptů \verb|faceid_server.py|,\newline \verb|extract_embeddings.py|a \verb|train.sh|.
\myparagraph{Rozšíření backendu o FaceID}
V prvním skriptu se nacházejí Tornado Handleři, jeden zajišťující komunikaci s frontendem v rámci posílání a ukládání fotek a druhý v rámci detekci obličeje ve fotce.
Předtím oba dva handleři ???fungovali??? v rámci jedné URL adresy. Já je rozdělil. Aplikace funguje následovně:
\begin{enumerate}
    \item Nepřihlášený uživatel se může buď přihlásit přes uživatelské jméno a heslo (vstupní stránka), přihlásit se přes FaceID nebo se zaregistrovat
    \item Po registraci musí uživatel čekat na schválení adminem, jinak se nepřhlásí (ani v případě, že již fotky má uložené v rámci biometické autentizace)
    \item Přihlášený uživatel má přístup k celé aplikaci, včetně možnosti nafotit si své fotky pro FaceID a natrénovat model pro autentizaci obličejů
\end{enumerate} 

\myparagraph{Uživatelské účty a identity}
Máme celkem 9 uživatelských účtů pro příhlášení, pro: 
\begin{itemize}
    \item 4 vyučující% Jan Švec, Vlasta Radová, Martin Bulín (bulinm) a Tomáš Lebeda (TODO)
    \item 3 členy týmu
    \item 1 admin účet, ke kterému mají přístup všichni členové týmu
\end{itemize}
Identit je 8. Jsou vytvořené pro všechny členy týmu, jeden kategorie `unknown' (aby se `nové' tváře nemohli přihlásit, jiank by je systém identifikoval jako validního uživatele) a 3 identity učitelé, aby mohlo dojít k právě zmíněné identity `unknown'.
Kategorie `unknown' má celkem 17 fotek, z toho 3 neobsahují lidský obličej, aby se neuronová síť naučila přiřazovat neplatné fotografie neautorizovanému uživateli.
Každý z členů má 10 až 25 fotek, které si každý sám nafotil.
Od každého vyučujícího jsme dostali 10 - 25 fotek. Každý s různým počtem, aby šlo testovat úspěšnost identifikace.
Rozhodovací práh je nastaven na alespoň 90\% pravděpodobnost identifikace uživatele. 
Je nutné podotknout, že identita `unkown' není vůbec uložená v databázi a nemá možnost se do aplikace přihlásit.
\myparagraph{Funkčnost přihlašování}
V případě, že se člověk přihlásí, je mu nastaven Json Web Token (JWT). Ten řeší zda je uživatelská session je validní či nikoliv.
Pokud jsou uživatelská credentials v pořádku a uživatel je zaregistován a schválený adminem je mu JWT uložen jako secure cookie.
Všechny další žádosti nyní automaticky kontrolují zda je JWT neexpiroval nebo s ním nebylo manipulovánu. 
JWT má totiž nastavenou donbu expirace jedné hodiny a obsahuje `podepsán' tajným klíčem, který je znám pouze serveru.  
\myparagraph{Odhlašování}
Pro odhlašování byl vytvořen Tornado Handler, který vymaže token obsahující cookie a přesměruje uživale zpět na přihlašovací stránku.

\section{Závěr}
Všechny položky zadání byly zodpovězeny.

\end{document}